% !TeX root = ../../tesis.tex
% §2.1 (entrada) — texto nuevo (#36)

Esta sección revisa cómo se ha estudiado una red de transporte representada
como grafo. Casi todos los trabajos son de 2015 en adelante, y una decena de
ellos propone o aplica indicadores de red; unos pocos anteriores se incluyen
porque fijaron los conceptos de base. El orden es por temas, no por autores.
Primero va la vulnerabilidad de las redes (§\ref{subsec:ea-vulnerabilidad}),
luego las herramientas y los datos abiertos con que se arma el grafo
(§\ref{subsec:ea-herramientas-datos}), después los indicadores
(§\ref{subsec:ea-indicadores}) y al final la topología anillo-radial con los
anillos de otras metrópolis (§\ref{subsec:ea-anillo-radial}). La síntesis
(§\ref{subsec:ea-sintesis}) junta tendencias, debates y vacíos, y ubica ahí el
problema de esta tesis.
